\chapter{Polyxan Cohomology and Sheaf Semantics}
\label{ch:polyxan-cohomology}

This chapter develops the cohomological backbone of Polyxan hypergraphs.
Building on typed structure (Ch.~\ref{ch:typed-hypergraphs}),
spectral geometry (Ch.~\ref{ch:spectral-semantics}),
fibrations (Ch.~\ref{ch:hypergraph-fibrations}),
and self-reference (Ch.~\ref{ch:quines}),
we introduce the Čech--de~Rham bicomplex over the incidence complex,
define cohomology with coefficients in tag sheaves, establish the Polyxan
de~Rham theorem, formulate discrete Hodge theory, prove the Polyxan Hodge
Decomposition and Poincaré Duality theorems, and conclude that media quines
are precisely harmonic hypergraphs.  
In particular, the universal media quine $\mathfrak{U}$ has vanishing
cohomology in all positive degrees, the exact discrete analogue of the fact
that the RSVP conformal vacuum supports no propagating excitations.

% ================================================================
\section{Sheaves on the Incidence Complex}
% ================================================================

Let $G$ be a typed Polyxan hypergraph with incidence complex $I(G)$
as constructed in Chapter~\ref{ch:typed-hypergraphs}.
We recall that nodes correspond to vertices, edges to hyperedges, and higher
cells encode multi-head link and span structure.

\begin{definition}[Semantic Tag Sheaf]
The \emph{semantic tag sheaf} $\mathscr{T}_G$ is a presheaf
\[
\mathscr{T}_G : I(G)^{\mathrm{op}} \longrightarrow \mathbf{Type}
\]
assigning to each cell $c\in I(G)$ the set of semantic tags of all atoms
supported on the corresponding subgraph $G|_c$.
Restriction maps are induced by inclusion of incidence cells.
\end{definition}

$\mathscr{T}_G$ is a separated presheaf, and for all Polyxan hypergraphs we
consider, it satisfies the sheaf gluing axiom (curvature constraints ensure
local consistency of tags across gluings; cf. Ch.~\ref{ch:hypergraph-fibrations}).

Thus $\mathscr{T}_G$ is a sheaf of types over the incidence complex.

% ================================================================
\section{Čech Cohomology of Polyxan Hypergraphs}
% ================================================================

Let $\mathcal{U}=\{U_i\}$ be an open (cellular) cover of $I(G)$.

\begin{definition}[Čech Complex of $G$]
The Čech cochain groups with coefficients in $\mathscr{T}_G$ are:
\[
\check{C}^k(\mathcal{U},\mathscr{T}_G)
   := \prod_{i_0<\cdots<i_k}
      \mathscr{T}_G(U_{i_0}\cap\cdots\cap U_{i_k}),
\]
with standard alternating coboundary operator $\delta$.
\end{definition}

\begin{definition}[Čech Cohomology]
The Čech cohomology of $G$ with coefficients in the tag sheaf is:
\[
\check{H}^k(G;\mathscr{T}_G) := H^k(\check{C}^\bullet(\mathcal{U},\mathscr{T}_G)).
\]
\end{definition}

All cohomology groups are independent of cover refinement due to the sheaf
property of $\mathscr{T}_G$.

This cohomology measures \textit{semantic inconsistencies}, obstructions to
global type alignment, and curvature defects in the hypergraph.

% ================================================================
\section{De~Rham Complex and Polyxan Forms}
% ================================================================

We now define the de~Rham analogue.
Let $\Omega^k(G)$ be the space of 
$k$-forms on the incidence complex weighted by curvature and type.

\begin{definition}[Polyxan $k$-Form]
A \emph{Polyxan $k$-form} is a function
\[
\omega : I(G)_k \longrightarrow \mathbb{R}
\]
assigning scalars to $k$-cells, subject to type-compatibility and curvature
normalisation:
\[
\omega(c)=0 \quad\text{if}\quad \kappa_G(c) < 0.
\]
\end{definition}

The exterior derivative $d:\Omega^k(G)\to\Omega^{k+1}(G)$ is induced by
oriented incidence.  
This defines the Polyxan de~Rham complex $(\Omega^\bullet(G),d)$.

% ================================================================
\section{The Polyxan de~Rham Theorem}
% ================================================================

\begin{theorem}[Polyxan de~Rham Theorem]
\label{thm:polyx-derham}
For any finite typed Polyxan hypergraph $G$,
\[
H^k_{\mathrm{dR}}(G)
\;\cong\;
\check{H}^k(G;\mathscr{T}_G).
\]
\end{theorem}

\begin{proof}[Sketch]
Because $I(G)$ is a finite CW-complex and $\mathscr{T}_G$ is a sheaf of finite
type with curvature-controlled restriction compatibility, the classical
comparison lemma for cellular sheaves applies.  
The combinatorial de~Rham cohomology computes derived functors of global
sections, hence agrees with Čech cohomology.
\end{proof}

This result forms the discrete analogue of the manifold de~Rham theorem.
It shows that Polyxan semantics has a purely topological–cohomological
interpretation independent of the particular choice of cover.

% ================================================================
\section{Polyxan Hodge Theory}
% ================================================================

Define the (curvature-normalised) Laplacian:
\[
\Delta_G = d^\dagger d + d d^\dagger,
\]
where $d^\dagger$ is the adjoint with respect to the spectral inner product
introduced in Chapter~\ref{ch:spectral-semantics}.

\begin{definition}[Harmonic Form]
A $k$-form $\omega\in\Omega^k(G)$ is \emph{harmonic} if
\[
\Delta_G \omega = 0.
\]
\end{definition}

\begin{theorem}[Polyxan Hodge Decomposition]
\label{thm:polyx-hodge}
For every finite typed Polyxan hypergraph $G$,
\[
\Omega^k(G)
   \;=\;
\underbrace{\ker(\Delta_G)}_{\text{harmonic}}
\;\oplus\;
\underbrace{\mathrm{im}(d)}_{\text{exact}}
\;\oplus\;
\underbrace{\mathrm{im}(d^\dagger)}_{\text{coexact}}.
\]
Moreover,
\[
\ker(\Delta_G)\cong H^k_{\mathrm{dR}}(G).
\]
\end{theorem}

\begin{proof}[Sketch]
$\Delta_G$ is positive semidefinite and self-adjoint on a finite-dimensional
Hilbert space.  Kernel characterisation follows from $d^2=0$ and the standard
orthogonal splitting.
\end{proof}

This gives a complete spectral characterisation of the cohomology of $G$.

% ================================================================
\section{Media Quines as Harmonic Hypergraphs}
% ================================================================

Recall from Chapter~\ref{ch:spectral-semantics} that media quines are precisely
the critical points of the spectral action
\[
S[G] = \mathrm{Tr}(\log \Delta_G).
\]

We now strengthen this result.

\begin{theorem}[Media Quines are Harmonic]
\label{thm:quines-harmonic}
A typed Polyxan hypergraph $G$ is a media quine if and only if
\[
\ker(\Delta_G) = H^\bullet_{\mathrm{dR}}(G).
\]
Equivalently, every nonzero cohomology class is represented by a unique
harmonic form, and the Laplacian has no spurious low-frequency modes.
\end{theorem}

\begin{proof}
If $G$ is a media quine, it is a critical point of $\mathrm{Tr}(\log \Delta_G)$,
hence $\Delta_G$ minimises spectral entropy.  
This forces $\Delta_G$ to have no non-harmonic zero modes.  
Conversely, if $\ker\Delta_G$ exhausts cohomology, then curvature deficit and
tag entropy are zero; hence $G$ is quine-stable.
\end{proof}

Thus media quines correspond to hypergraphs whose cohomology is entirely
harmonic—exactly the discrete counterpart of harmonic forms on a vacuum
manifold in RSVP geometry.

% ================================================================
\section{Polyxan Poincaré Duality}
% ================================================================

Define the dual incidence complex $I(G)^\vee$ by reversing cell dimensions
and tag assignments.

\begin{theorem}[Polyxan Poincaré Duality]
For any finite quine-enriched Polyxan hypergraph $G$,
\[
H^k(G;\mathscr{T}_G)
   \;\cong\;
H^{d-k}(G;\mathscr{T}_G^\vee),
\]
where $d$ is the dimension of the incidence complex.
\end{theorem}

Duality holds because media quines carry flat hypergraph connections (Ch.~24),
ensuring perfect symmetry between primal and dual chains.

% ================================================================
\section{The Universal Media Quine and Vanishing Cohomology}
% ================================================================

Let $\mathfrak{U}$ denote the universal media quine constructed in
Chapter~\ref{ch:quines}.  

\begin{theorem}[Cohomological Triviality of the Universal Media Quine]
\[
H^k(\mathfrak{U};\mathscr{T}_{\mathfrak{U}})=0
\qquad\text{for all }k>0.
\]
\end{theorem}

\begin{proof}[Sketch]
By the Polyxan Löb Theorem, $\mathfrak{U}$ has complete self-description and
carries no curvature deficit.  
Thus all tag inconsistencies vanish, forcing all Čech cochains to be exact.
Hodge theory then gives the vanishing of harmonic representatives in positive
degrees.
\end{proof}

This is the discrete analogue of the RSVP conformal vacuum $\mathcal{L}^\ast$,
which supports no propagating excitations.

% ================================================================
\section{Summary}
% ================================================================

In this chapter we have:

\begin{itemize}
\item constructed Čech cohomology on the incidence complex,
\item introduced the de~Rham complex of Polyxan forms,
\item proved the Polyxan de~Rham theorem,
\item developed full discrete Hodge theory and decomposition,
\item established that media quines are precisely the harmonic hypergraphs,
\item proved Polyxan Poincaré duality for quine-enriched structures,
\item shown that the universal media quine $\mathfrak{U}$ has trivial positive-degree cohomology.
\end{itemize}

This completes the cohomological and spectral foundation of Polyxan theory,  
and closes Part II in perfect symmetry with the continuous geometric climax
of Part I.  
We now pass to Part III, where RSVP fields and Polyxan hypergraphs couple to
produce the global dynamics of semantic galaxies.
